\section*{600 К. Две фазы и сходимость контактов}
Низкотемпературный переход даёт два хорошо разделённых максимума
распределения порядка. С увеличением ячейки барьер между ними растёт,
а условные средние составы в бассейнах меняются значительно слабее.
В квадратных скобках приведены 95\%-е блочные интервалы.
Столбец \(\beta\Delta F\) даёт высоту барьера в профиле
\(-\ln P(Q)\) при равных высотах пиков. Сам состав определён по
равенству статистических весов отдельных вариантов фаз.

% CONTROL_TABLE

На 1400 узлах 32 зонтичных окна дали поле
\(\Delta\mu^*=3{,}156552\)~эВ, 95\%-й интервал [3,156251; 3,156886].
Средний порядок составляет 0,030 в \(\alpha\) и 0,970 в \(\alpha_2\),
барьер достиг \(35{,}7\,k_{\rm B}T\). Изменение контактов относительно
768 атомов составляет \(-0{,}005\) и \(-0{,}011\) п.п. Оценки числа
эффективных кадров с поправкой на корреляции --- 47 и 136 соответственно.
Восемь независимых траекторий без зонтика при полях 3,156 и 3,159~эВ
дали после приведения к \(\Delta\mu^*\) 5,81--5,83 и 24,30~ат.\% Al.
Разности с основным расчётом не превышают 0,06 п.п. Блочные интервалы
этих проверок и зонтичного расчёта перекрываются.

Ужесточение ограничения до \(W<0{,}35\) на 1400 атомах изменяет
\(x_\alpha\) на \(+0{,}006\)~п.п. (95\% от \(-0{,}104\) до \(+0{,}135\)),
а \(x_{\alpha_2}\) на \(+0{,}048\)~п.п. (от \(+0{,}017\) до \(+0{,}087\)).
Небольшая остаточная чувствительность правой границы различима
при достигнутой статистике. Эти две стенки характеризуют конечное
изменение ограничения в проверенном диапазоне.

\subsection*{Откуда берётся расхождение при 600 К}
В базе Liang ширина равновесного интервала составляет 9,11~п.п.,
в прямом MTP --- 18,40~п.п. Левая граница MTP ниже оценки Liang на 5,06~п.п.,
правая выше на 4,23~п.п. В упорядоченной фазе Al занимает 97,05\% мест
своей номинальной подрешётки. Для сравнения, при
1000~К прямой MTP и CALPHAD дают близкие левые границы
(14,98 и 14,75), тогда как правая граница прямого расчёта уже там
лежит на 1,44~п.п. выше значения Liang.

Когда форма ячейки свободна, \(\langle c/a\rangle\) при 600~К составляет
1,57764 в \(\alpha\) и 1,61012 в \(\alpha_2\). При фиксированном
авторском отношении осей на 768 атомах прямой MTP даёт
\(\FixedAlpha/\FixedOrdered\)~ат.\% Al. Для тех же механических условий
классический C46 даёт 4,55/24,10 на 5120 узлах. При согласованном
\(c/a\) прямой расчёт сдвигает левую границу относительно C46 на
\(+3{,}03\) п.п., правую --- на \(-0{,}29\) п.п.

При сравнении C46 с прямым MTP при фиксированном \(c/a\) обе модели
используют классическую статистику и один потенциал. Остаточная
разность включает аппроксимацию свободной энергии кластерным базисом,
квазигармоническое приближение и эффект ограничения атомных смещений
в прямом расчёте. Разделить эти вклады по двум равновесным составам
нельзя. Для оценки ошибки самого MTP при 600~К требуется независимый
ориентир, особенно для растворимости Al в \(\alpha\).
